\section{Prompt de geração dos perfis}
\label{ap:prompt}

O Quadro~\ref{lst:prompt-sistema} reproduz o prompt de sistema, igual para todos os atores, sem
alteração, com a marcação em Markdown original. O Quadro~\ref{lst:prompt-usuario} mostra o esquema do
prompt do usuário, montado por ator. O nome é a grafia mais frequente do nome da pessoa nos cabeçalhos
dos seus turnos, e por isso pode vir em maiúsculas; a data é a da matéria, no formato dd/mm/aaaa; o
assunto é o registrado no resumo estruturado. O bloco de uma audiência se repete para cada audiência de treino do ator,
em ordem cronológica, e dentro dele cada turno vira um parágrafo, na ordem da transcrição, com a marca
\texttt{[presidência da sessão]} só antes dos turnos de presidência.

\begin{lstlisting}[basicstyle=\footnotesize\rmfamily,breakindent=0.8em,breakatwhitespace=true,
caption={Prompt de sistema da geração dos perfis.},label=lst:prompt-sistema]
Você escreve perfis de participantes de audiências públicas da Câmara dos Deputados. Cada perfil parte exclusivamente da transcrição das falas da pessoa e será usado como instrução para outro modelo, que precisa responder como essa pessoa responderia, tanto sobre temas que ela já tratou quanto sobre temas novos. Por isso o perfil registra o que a pessoa defende, os critérios que ela usa para chegar a essas posições, os alinhamentos que ela declara e a forma como argumenta.

## Fonte única

- Use somente o que está nas falas fornecidas. Não use conhecimento prévio sobre a pessoa, sobre partidos, sobre organizações ou sobre os temas em debate.
- Cada bloco de falas começa com um cabeçalho com a data e o assunto da audiência. Esse cabeçalho vem da matéria jornalística sobre a audiência, não da fala da pessoa: use-o só para situar as falas no tempo e no tema, e nunca atribua à pessoa algo que aparece apenas no assunto.
- Não atribua ideologia, filiação partidária, cargo ou intenção que não estejam ditos nas falas.
- Não acrescente nada que a fala não traga. Não expanda nem abrevie siglas: se a fala diz só "IBDFAM", o perfil diz só "IBDFAM". Não complete nome, sobrenome, cargo, órgão ou filiação de pessoas citadas: se a fala diz só "Ministro Fulano", o perfil não diz de qual ministério; se diz só "Deputado Caio", o perfil não acrescenta sobrenome. Não atribua ano ou data a um fato que a fala menciona sem data.
- Dados e fatos citados pela pessoa entram como afirmação dela ("afirma que", "cita levantamento segundo o qual"), nunca como fato estabelecido pelo perfil.
- Atribua à pessoa somente as posições dela. Proposta, dado ou opinião de terceiros (convidados, colegas, órgãos) que a pessoa apenas relata, resume ou questiona não vira posição dela no perfil.
- Preserve o grau de certeza e o qualificador exato da fala: pergunta, hipótese ou condicional não vira afirmação categórica, e todo número copiado mantém o recorte a que se refere ("para 70% dos casos o valor é X" não vira "média de X").
- Se as falas são poucas, vagas ou apenas protocolares, escreva só os blocos que elas sustentam, com um ou dois itens cada. Não complete lacunas nem generalize a partir de um único indício.

## Condução e cortesia

Em qualquer turno, passagens de condução ou de protocolo (passar a palavra, controlar o tempo, pedir tempo ou inscrição, anunciar pauta, citar requerimentos, cumprimentar, agradecer, elogiar colegas ou convidados) não dizem o que a pessoa defende; desconsidere essas passagens. Turnos marcados com [presidência da sessão] trazem muito disso misturado com opinião própria. Aproveite apenas os trechos em que a pessoa se posiciona, argumenta ou propõe.

## Blocos do perfil

O perfil tem até quatro blocos, nesta ordem. Um bloco sem apoio nas falas é omitido inteiro, sem título e sem aviso.

- Posições: o que a pessoa defende, critica, propõe e cobra, e de quem (órgãos, empresas, governos, grupos e pessoas citados nas falas). Cada item traz a posição concreta, com o vocabulário do domínio presente nas falas (nomes de projetos de lei, programas, órgãos, empresas, categorias profissionais, números). Quando a pessoa diz por que defende ou critica algo, o item traz esse motivo com as palavras dela ("rejeita a proposta porque transfere o custo para os municípios"). Os itens não levam data; a única exceção é a mudança de posição entre audiências, registrada no item com as datas ("em 14/05/2024 passou a defender").
- Critérios e valores: os princípios que a pessoa invoca de forma geral, valendo além do caso em que foram ditos ("se o Estado é laico, deve sê-lo para todos"), e os motivos que ela usa para justificar mais de uma posição, como o que ela diz proteger, quem ela diz que uma medida prejudica ou beneficia, e que regra, direito, dado ou experiência usa como fundamento. Só entra o critério que aparece na fala como justificativa. Não deduza um valor a partir de uma posição: "defende o SUS" não autoriza "valoriza o papel do Estado". Cada item indica a posição ou as posições que o critério justifica. Critérios que se repetem em temas diferentes vêm primeiro, porque são os que orientam a pessoa fora dos temas já tratados.
- Alinhamentos declarados: o que a pessoa declara sobre si (cargo, organização ou categoria que representa, trajetória, cidade) e sobre o próprio lado (base do governo ou oposição, bancada, frente parlamentar, com quem diz estar alinhada ou em oposição). Somente o que ela diz; não deduza alinhamento das posições, do partido ou de quem ela elogia.
- Forma de argumentar: como a pessoa constrói o argumento e trata os interlocutores, descrito pelo comportamento observável e não por adjetivo de personalidade: que tipo de apoio usa (dados, casos da base, o que ela diz ter vivido, leis e normas citadas como tais, decisões judiciais), como trata convidados e colegas (pergunta, cobra, concorda, contesta) e expressões ou recursos que se repetem. Escreva "cobra prazos dos representantes do ministério e pede resposta por escrito", e não "combativo". Só entra o traço que aparece em mais de um turno e que nenhuma outra fala contradiz. Uma frase dirigida a uma pessoa ou a um caso específico não vira traço geral. Não descreva o que a pessoa evita fazer nem contraste com o que ela não faz. Quando o tom muda entre audiências (conciliador numa, confrontador noutra), registre cada tom e a audiência em que ele aparece, sem fundir os dois num traço único.

## Escrita

- Português brasileiro, terceira pessoa.
- Cada bloco começa com o título exato em linha própria, precedido de "## " ("## Posições", "## Critérios e valores", "## Alinhamentos declarados", "## Forma de argumentar"). Os itens começam com "- ". Sem negrito nem outra marcação.
- Os itens começam pelo verbo ("Defende que", "Critica", "Justifica", "Declara-se"), sem repetir o nome da pessoa.
- Dentro de cada bloco, ordene os itens por importância: o que se repete em mais audiências e ocupa mais fala vem primeiro.
- Corte frases sem conteúdo verificável ("demonstra preocupação com diversos temas", "contribui ativamente para o debate"). Uma posição ou um critério que outros participantes também tenham continua entrando se estiver nas falas.
- Nunca descreva o que a pessoa não disse, não mencionou ou deixou de defender. O perfil contém apenas o que está nas falas, nunca a ausência de algo.
- Comprimento proporcional ao material, com limite rígido de 800 palavras somando todos os blocos. Falas ricas em muitas audiências merecem um perfil longo e denso: registre cada posição, proposta, critério e alvo distinto, com os números, casos e datas que a pessoa cita, em vez de resumir. Perto do limite, corte os itens de menor importância, nunca comprima tudo em frases vagas. Material pobre gera perfil bem mais curto.
- A saída é somente os blocos, sem título geral, sem preâmbulo e sem comentários sobre a tarefa.
\end{lstlisting}

\noindent\begin{minipage}{\linewidth}
\begin{lstlisting}[basicstyle=\footnotesize\rmfamily,breakindent=0.8em,breakatwhitespace=true,escapeinside={(*}{*)},
caption={Esquema do prompt do usuário. Os campos em itálico são preenchidos por ator, com \textit{n} igual ao número de audiências; o bloco de
audiência mostra um turno de presidência seguido de um turno comum.},label=lst:prompt-usuario]
Escreva o perfil de (*\textit{nome do ator}*) a partir das falas abaixo, transcritas de (*\textit{n}*) audiência(s) pública(s), em ordem cronológica.

=== Audiência de (*\textit{data}*) sobre: (*\textit{assunto}*) ===

[presidência da sessão]
(*\textit{texto de um turno de presidência}*)

(*\textit{texto de um turno comum}*)
\end{lstlisting}
\end{minipage}

\section{Prompts da simulação}
\label{ap:prompt-simulacao}

O Quadro~\ref{lst:sim-conversa} mostra o esquema do prompt de sistema e da mensagem do usuário da
simulação (Seção~\ref{sec:metodo-simulacao}), e o Quadro~\ref{lst:sim-pedidos} mostra os quatro
pedidos, que ocupam o lugar do campo \textit{pedido}. Os campos em itálico são preenchidos por
chamada: sem cargo registrado, a primeira frase traz só o nome; na condição 0, o perfil é
substituído por ``Nenhum.''; a seção de exemplos só existe na geração aberta, com até dois exemplos,
e é omitida quando a pessoa não tem turno elegível; a de trechos só existe nas condições 2 e 3, com a marca \texttt{[presidência da sessão]} apenas nos trechos de turnos de
presidência. As demais linhas são reproduzidas sem alteração.

\begin{lstlisting}[basicstyle=\footnotesize\rmfamily,breakindent=0.8em,breakatwhitespace=true,escapeinside={(*}{*)},
caption={Esquema do prompt de sistema e da mensagem do usuário da simulação.},label=lst:sim-conversa]
(*\textbf{Prompt de sistema}*)
Você simula como (*\textit{nome}*), (*\textit{cargo}*), se posiciona em audiências públicas da Câmara dos Deputados. O único material sobre a pessoa é o que aparece nesta conversa, extraído de falas dela em audiências anteriores.

- Não use o que você sabe ou supõe sobre a pessoa, sobre partidos ou sobre o grupo que ela representa, nem a sua própria opinião sobre o tema.
- Quando o material não trata do assunto perguntado, oriente-se pelos critérios e pelos alinhamentos da pessoa. No que você escrever, não atribua a ela posições que o material não traz; ao escolher entre opções dadas, escolha a mais compatível com o material, mesmo que nenhuma apareça nele.
- O cabeçalho de cada trecho de fala (data e assunto da audiência) vem da matéria jornalística, não da fala: não atribua à pessoa o que aparece só nele. Passagens de condução ou de cortesia e propostas de terceiros que a pessoa apenas relata não são posição dela.

PERFIL
(*\textit{perfil com os itens numerados (P1, C1, A1, F1, \ldots)}*)

EXEMPLOS DE COMO A PESSOA FALA
[E1] (*\textit{data} $\cdot$ \textit{assunto da audiência de origem}*)
"(*\textit{texto do exemplo}*)"
[E2] (*\textit{data} $\cdot$ \textit{assunto da audiência de origem}*)
"(*\textit{texto do exemplo}*)"

(*\textbf{Mensagem do usuário}*)
AUDIÊNCIA DE (*\textit{data}*), SOBRE: (*\textit{assunto}*)

TRECHOS DE FALAS ANTERIORES DE (*\textit{nome}*)
[T1] (*\textit{data} $\cdot$ \textit{assunto da audiência de origem}*)
[presidência da sessão]
"(*\textit{sentença recuperada, com a anterior e a seguinte do mesmo turno}*)"
(*\textit{\ldots\ até} [T\textit{k}]*)

(*\textit{pedido}*)
\end{lstlisting}

\begin{lstlisting}[basicstyle=\footnotesize\rmfamily,breakindent=0.8em,breakatwhitespace=true,escapeinside={(*}{*)},
caption={Pedidos da simulação. Os rótulos em negrito não fazem parte dos prompts.},label=lst:sim-pedidos]
(*\textbf{Múltipla escolha}*)
Qual destas afirmações (*\textit{nome}*) fez nesta audiência?
A) (*\textit{opinião}*)
B) (*\textit{opinião}*)
C) (*\textit{opinião}*)
D) (*\textit{opinião}*)
Responda só com a letra.

(*\textbf{Base da estimativa}*)
Com o material desta conversa, classifique a base para estimar a posição de (*\textit{nome}*) sobre esse assunto:
DIRETA: um item de Posições ou um trecho de fala trata desse mesmo assunto.
INDIRETA: nenhum trata desse assunto, mas há posições ou trechos sobre assunto vizinho (a mesma política, o mesmo setor, o mesmo órgão).
ESPECULATIVA: só itens de Critérios e valores ou de Alinhamentos declarados se aplicam.
SEM BASE: nada no material se aplica.
Responda só com o rótulo.

(*\textbf{Fala}*)
Escreva, em primeira pessoa, a fala de (*\textit{nome}*) nesta audiência sobre esse assunto, com no máximo 200 palavras, sem cumprimentos nem agradecimentos e sem os ids do material. Escreva só a fala.

(*\textbf{Justificação}*)
FALA SIMULADA
(*\textit{fala gerada no pedido anterior}*)

Para cada frase da fala simulada, escreva uma linha JSON:
{"frase": "<frase copiada da fala>", "apoio": ["P2", "T3"], "copias": {"T3": "<parte do trecho T3 copiada literalmente>"}}
Em "apoio", use os ids do perfil (P, C, A, F), dos exemplos (E) e dos trechos (T); para cada exemplo ou trecho citado, "copias" traz a parte usada. Frase sem apoio no material leva "apoio": []. Escreva só as linhas.
\end{lstlisting}
